% PDF 217--224；原书第 189--196 页。
\section{基于人类反馈的强化学习的演变}
\label{sec:ch05-evolution}
下面考察几项重要研究，说明它们怎样推动使用强化学习改善大语言模型生成文本的质量，使模型在对话中产生更接近人类预期的结果。

\subsection{大语言模型的安全性、质量与有据性}
评估生成模型并不容易，尤其是生成开放文本的对话模型。若模型有明确的目标，可以直接比较预测和标签。但如果输出没有确切答案，例如聊天机器人与用户对话，就很难确定怎样用数学方法衡量结果的质量。LaMDA 系统在这个方向上引入了新的指标，对大语言模型的价值观对齐作出了重要贡献（\refcite{ch05-thoppilan}{Thoppilan et al., 2022}）。

\begin{keypoint}
\noindent\textbf{主要贡献}\quad LaMDA 系统为开放式对话系统的评估引入了趣味性、安全性、有据性、信息性等新指标。这些指标补充了已有的合理性和具体性准则，完善了对话系统在质量、安全性和有据性方面的评估指标。
\end{keypoint}

LaMDA 是 Google 开发与维护的一系列为文本生成优化的语言模型。评估 LaMDA 时采用三类基本指标，即质量（quality）、安全性（safety）和有据性（groundedness）。以这些指标衡量 LaMDA 的表现和效果，可从多个方面评估模型生成优质、安全、具有事实根据的对话的能力。下面说明这些目标及评估 LaMDA 表现所用的指标。

\begin{itemize}
\item \textbf{质量}是第一项目标，包括合理性、具体性和趣味性（sensibleness, specificity, and interestingness，SSI）三个方面，由人工评分员评估。
  \begin{itemize}
  \item 合理性评估模型的回答在对话语境中是否连贯，是否避免了常识错误、荒谬说法及自相矛盾。
  \item 具体性衡量回答是否针对具体的对话语境，是否避免泛泛或含糊的陈述。
  \item 趣味性评估模型生成富有见解、出人意料或机智的回答的能力，这些特征有助于提高对话质量。
  \end{itemize}
\item \textbf{安全性}是第二项目标，与负责任地开发和部署人工智能有关。安全性指标由一组规则组成，规定对话中希望模型采取的行为，以防止意外后果、用户遭受风险和不公平的偏见。这些目标引导模型避免生成含有暴力或血腥内容的回答，避免助长侮辱性言辞或仇恨性的刻板印象，也避免粗言秽语。
\item \textbf{有据性}是第三项目标，处理语言模型生成看似合理却与事实不符的陈述的问题。它衡量模型关于外部世界的回答中，有权威外部来源可以证实的回答所占的比例。与之相关的指标是信息性（informativeness），量化有已知来源提供信息支持的回答比例。使 LaMDA 的回答依据已知来源，本身无法保证事实准确，但用户或外部系统可以根据这些来源的可靠程度判断回答是否有效。
\end{itemize}

\subsubsection{方法}
对话模型 LaMDA 的训练分为预训练和微调两个阶段。

\begin{enumerate}
\item \textbf{LaMDA 预训练}

预训练阶段从公开的对话数据和网页文档中汇集了 1.56 万亿个词，数据规模超过此前的模型。经词元切分得到 2.81 万亿个 SentencePiece 词元后，用 GSPMD（\refcite{ch05-xu}{Xu et al., 2021}）训练 LaMDA，使其根据前文预测后续词元。预训练的 LaMDA 模型可用于自然语言处理的许多研究方向，包括程序合成、零样本学习、风格迁移以及参加 BIG-bench 研讨会。

\item \textbf{LaMDA 微调}

在微调阶段，训练 LaMDA 同时执行生成和分类任务，如图~\ref{fig:ch05-lamda-classifiers} 所示。在生成方面，它根据指定上下文生成自然语言回答。在分类方面，它对回答的安全性和质量（SSI）评分。由此得到能够执行这两种功能的统一多任务模型。训练 LaMDA 生成器时，使用由两人来回交谈组成的对话数据集，预测下一个词元。同时，用带标注的数据训练 LaMDA 分类器，使其依据各自的上下文对回答的安全性和质量（SSI）评分。

在对话交互中，LaMDA 生成器根据正在进行的多轮对话，生成多个候选回答。然后，LaMDA 分类器为每个候选回答给出 SSI 和安全性分数。安全性分数较低的回答会被舍弃，其余候选按 SSI 分数重新排序，得分最高的作为最终输出。为了提高候选回答的质量，生成任务所用的训练数据还经过 LaMDA 分类器的额外筛选，以增加高质量候选的比例。
\end{enumerate}

\begin{figure}[!htbp]
\centering
% 重绘自原书第 191 页的图 5.4，保留各候选字符串和分数。
\begin{tikzpicture}[
  font=\footnotesize\sffamily,
  chip/.style={draw=black!30,fill=black!6,align=center,rounded corners=2pt,
    inner sep=2.3pt,minimum height=5mm},
  flow/.style={draw=black!60,-{Stealth[length=1.25mm]},line width=.50pt}]
\node[chip] (hi) at (0,0) {Hi};
\node[chip,fill=blue!19,minimum width=15mm] (gen) at (1.38,0) {生成器};
\node[chip,minimum width=28mm,font=\scriptsize\ttfamily] (first) at (4.03,.53) {"Hi! How's it going?"};
\node[chip,minimum width=28mm,font=\scriptsize\ttfamily] (second) at (4.03,0) {"Hello! How are you?"};
\node[chip,minimum width=28mm,font=\scriptsize\ttfamily] (third) at (4.03,-.53) {"Greetings!"};
\node[chip,fill=red!19,minimum width=16mm] (cls) at (6.80,0) {分类器};
\node[chip,fill=green!18,align=left,inner sep=3pt,font=\scriptsize\sffamily] (score) at (9.20,0)
  {安全性 90\%\\合理性 80\%\\具体性 70\%\\趣味性 50\%};
\draw[flow] (hi.east)--(gen.west);
\draw[flow] (gen.east)--(2.47,0);
\draw[flow] (5.59,0)--(cls.west);
\draw[flow] (cls.east)--(score.west);
\draw[black!65] (2.64,.82)--(2.46,.82)--(2.46,-.82)--(2.64,-.82)
 (5.42,.82)--(5.60,.82)--(5.60,-.82)--(5.42,-.82);
\end{tikzpicture}

\caption{LaMDA 通过生成与分类得到各项分数。}
\label{fig:ch05-lamda-classifiers}
\end{figure}

这项研究还考察了怎样提高 LaMDA 回答的事实准确性。研究者整理了一个数据集，加强 LaMDA 初次回答所用的知识基础。数据集由人与 LaMDA 的对话组成，在适用时加入信息检索查询和对应的搜索结果。随后，在这个数据集上微调 LaMDA 的生成器与分类器，教模型在与用户交互时使用外部信息检索系统，改善回答所依赖的事实根据。自 LaMDA 最初开发以来，从外部检索信息在大语言模型应用中已越来越常见。

图~\ref{fig:ch05-lamda-grounding} 说明 LaMDA 怎样与外部信息检索系统交互，使回答有据可查。如前所述，过程从调用 LaMDA-Base 模型开始，即图中的蓝色组件。用户的查询或提示构成模型输入，以黄色显示。LaMDA-Base 根据输入生成一个回答。随后调用 LaMDA-Research 模型继续处理。通过依次调用，形成逐步改进的迭代过程。LaMDA-Research 输出的第一个词决定查询信息检索系统还是直接回答用户。这个词是一个标识符，指明下一步的接收方。LaMDA-Research 发现需要补充信息，或要加强回答的事实根据时，就向外部信息检索系统发起查询。信息检索系统根据查询生成响应，图中以绿色显示，再将其纳入正在进行的对话。

\begin{figure}[!htbp]
\centering
\includegraphics[trim=0 156bp 0 159bp,clip,width=\linewidth]{assets/fig05_05.pdf}
\caption{LaMDA 通过与外部信息检索系统交互，使回答有据可查。蓝色为 LaMDA 模型，黄色为送入模型的输入，红色为 LaMDA 生成的响应，绿色为信息检索工具集（TS）的响应。示例查询埃菲尔铁塔何时建成，工具先后返回 1887 年 1 月 28 日开工、1889 年 3 月 31 日开放，模型据此答复用户。}
\label{fig:ch05-lamda-grounding}
\end{figure}

\subsubsection{评估与结果}
评估时，从两人参与的多轮对话中收集预训练模型、微调模型以及人给出的回答。然后，人工评分员按质量、安全性和有据性指标评估。结果表明，在各个维度和不同模型规模下，LaMDA 均优于预训练模型，如图~\ref{fig:ch05-lamda-scores} 所示。质量指标随模型参数量增大而改善。仅扩大模型规模并未改善安全性，微调则能改善这一点。有据性与模型规模正相关，而微调使模型更容易利用外部知识来源。

\sourcefigure{fig05_06.pdf}{比较预训练模型（PT）、微调模型（LaMDA）及人类评分员生成的对话（Human）。比较指标为合理性、具体性、趣味性、安全性、有据性和信息性。}{fig:ch05-lamda-scores}

\subsection{基于人类偏好的深度强化学习}
开发安全的人工智能系统，一个重要方向是摆脱人工编写的目标函数。若用过于简单的替代目标，或所设目标与复杂目标稍有偏离，人工智能也可能产生有害甚至危险的行为。为解决这个问题，OpenAI 与 DeepMind 的安全团队合作提出了一种算法。通常会设计奖励函数，以取得来自环境的奖励。这个研究则让人类监督者发挥关键作用，说明自己对不同选项（轨迹片段）的“偏好”（\refcite{ch05-christiano}{Christiano et al., 2017}）。

\begin{keypoint}
\noindent\textbf{主要贡献}\quad 这项研究表明，即使缺乏专业技术知识，人也能有效指导强化学习系统。这种人工智能系统通过反复试错来学习。人无需为算法预先定义目标，强化学习系统可以直接从人的指导中学到复杂的目标。
\end{keypoint}

\subsubsection{方法}
这种方法根据人类偏好调整奖励函数，同时训练一个最大化预测奖励的策略，以符合人类偏好。研究人员让人比较展示智能体行为的短视频片段，无需为它们打分。这样，人比较起来更容易，同时仍能有效学习偏好。比较短视频很快，也比比较单个状态能提供更多信息。研究还指出，收集实时反馈可提高系统表现，并防止系统利用学得的奖励函数中的弱点。

论文中的训练过程围绕一个三步反馈循环展开，涉及人的反馈、智能体对目标的理解，以及强化学习训练，如图~\ref{fig:ch05-reward-predictor} 所示。

\begin{enumerate}
\item \textbf{第 1 步。}智能体通过多步操作与环境交互，在每一步收到观测 $O_t$，并采取动作 $A_t$。传统做法由环境提供奖励 $r_t$，引导智能体最大化累计奖励。在这项研究中，作者假定有一位人类监督者，可以说明自己对不同轨迹片段的“偏好”。为此，引入可学习的策略 $\pi\colon O\to A$ 和奖励函数估计 $\hat r\colon O\times A\to\mathbb{R}$，二者均以深度神经网络参数化。策略 $\pi$ 与环境交互，生成一组轨迹 $\tau_1,\ldots,\tau_2$。再用传统强化学习算法更新策略参数，以最大化预测奖励 $r_t=\hat r(o_t,a_t)$ 的总和。
\item \textbf{第 2 步。}从生成的轨迹 $\tau_1,\ldots,\tau_2$ 中选取成对的轨迹片段 $\sigma_1,\sigma_2$，交给人类监督者比较。监督者对这些片段进行评估，并给出相对偏好的反馈。
\item \textbf{第 3 步。}用监督学习优化奖励函数估计 $\hat r$ 的参数，使估计结果符合截至此时从人类监督者收集到的偏好。
\end{enumerate}

\begin{figure}[!htbp]
\centering
% 重绘自原书第 195 页的图 5.7。
\begin{tikzpicture}[
  font=\footnotesize\sffamily,
  block/.style={draw=cyan!29!black!45,fill=cyan!19,rounded corners=2.6pt,align=center,
    minimum width=26mm,minimum height=12mm,inner sep=3pt},
  wire/.style={draw=blue!25!black!66,-{Stealth[length=1.4mm]},line width=.52pt}]
\node[block] (algo) at (0,0) {强化学习算法};
\node[block] (env) at (7.33,0) {环境};
\node[block,fill=yellow!30,minimum height=8mm,minimum width=31mm] (predict) at (3.66,2.36) {奖励预测器};
\node[align=center] (hum) at (7.04,2.36) {人类反馈};
\node[align=center] at (.0,2.33) {预测奖励};
\draw[wire] (predict.west) to[out=180,in=78] (algo.north);
\draw[wire,densely dotted] (hum.west) -- (predict.east);
\draw[wire] ([yshift=2.9pt]env.west) to[out=180,in=-90] (predict.south);
\draw[wire] ([yshift=2.9pt]env.west) -- node[below] {观测} ([yshift=2.9pt]algo.east);
\draw[wire] ([yshift=-3.4pt]algo.east) -- node[below] {动作} ([yshift=-3.4pt]env.west);
\end{tikzpicture}

\caption{训练奖励预测器与比较轨迹片段是异步的。智能体的目标是使预测奖励最大。}
\label{fig:ch05-reward-predictor}
\end{figure}

策略 $\pi$、奖励函数估计 $\hat r$ 和人类反馈的处理流程异步运行，按第 $1\to 2\to 3\to 1$ 步依次循环。

在优化算法方面，作者选择了在奖励函数改变时仍有效的一类策略优化算法，即策略梯度方法（policy gradient methods）。这些方法能有效更新策略 $\pi$，包括用于 Atari 游戏的优势行动者—评论家算法（Advantage Actor Critic），以及用于 MuJoCo 仿真的信赖域策略优化（trust region policy optimization）。

在人类反馈的处理流程中，从策略抽取两个轨迹片段，以 1 至 2 秒的短视频展示给人类监督者。监督者可以选更喜欢的一项，也可以认为两项同样可取，或者两项都不可取。数据库 $D$ 记录轨迹片段对 $(\sigma_1,\sigma_2)$，以及在 1、2 上的均匀分布 $\mu$。原书写作，偏好 $\sigma_1$ 时 $\mu$ 为 1，偏好 $\sigma_2$ 时为 2，两项同样可取时为 1.5。两项都不可取的轨迹对不计入数据库 $D$。\footnote{译注：Christiano 等人的论文第 2.2.2 节把 $\mu$ 定义为 $\{1,2\}$ 上的概率分布。若只偏好一项，概率全在该项上。两项同样可取时，$\mu$ 为均匀分布。无法比较时不记录。见 \url{https://arxiv.org/abs/1706.03741}。}

拟合奖励函数时，训练模型从收集到的轨迹偏好中推断奖励函数。作者把偏好建模为由 Bradley–Terry 模型（或 Boltzmann 理性模型）生成。偏好轨迹 A 胜过轨迹 B 的概率与两条轨迹回报之差的指数值成比例。这样，回报之差就可以作为二分类问题的 logit。因此，使用交叉熵损失来训练奖励函数，使其准确预测偏好比较的结果。

\subsubsection{评估与结果}
研究者在颇具挑战的 Atari 游戏 Enduro 中评估系统表现。由于游戏复杂且试错学习有限，传统强化学习网络难以处理这个游戏。纳入人类反馈后，系统最终在 Enduro 中达到了超过人类玩家的结果，表明了人类指导的作用。研究还发现，系统在其他游戏及机器人仿真任务中的表现，与标准强化学习设置相当，但在 Qbert 和 Breakout 等游戏中遇到了困难，甚至失败。

此外，研究者还考察系统能否理解并执行人指定的目标，即便环境没有给出明确提示。训练智能体学习多种新行为，包括后空翻、单腿行走，以及在 Enduro 中与另一辆车并排行驶。最后一项有别于通过超车使游戏分数最大化的常规目标。实验旨在评估系统能否学习和执行超出游戏环境所定义的即时目标的特定行为。
